% STATUS: written 2026-11-05 (every class row points to a section read for this
% chapter; excerpts in ledger/projects/derivations-encyclopaedia/excerpts/).
% Moved 2026-11-07 from chapters/21-open-problems/02-by-class into Chapter 2,
% with the problems themselves now stated in their own chapters.
% Sources: Freudenburg 2017, Biswas--Gurjar--Kolte 2014, Belov-Kanel--Kontsevich
% 2007, Hadwin--Li 2004, Ayupov--Kudaybergenov 2012/2015,
% Ayupov--Arzikulov 2014, Ayupov--Yusupov 2020, Habibzadeh--Sahleh 2019,
% Bresar--Villena 2002, Bresar 1993 TAMS, Bresar 2005, Martindale 1964,
% Benkovic 2014, Vukman 1990, Lee 1997.
\section{Open problems by class of algebra}\label{sec:basic-theory:open-problems-by-class}

The chapters state each problem in the section that asks its question. % claim-ok: own description of the arrangement of this section
This section keeps the same problems and sorts them by the algebra they live
in, so that a reader who works on one class sees the frontier of that class.
For each class it records which problem is settled there, which is not, and the
chapter that carries the discussion.

The first class is the polynomial algebra $k[x_1,\dots,x_n]$ over a field of
characteristic $0$, and it holds the deepest problems of the book.
Chapter~\ref{ch:commutative} reads them in
Section~\ref{sec:commutative:zariski-lipman} and
Section~\ref{sec:commutative:jacobian}. % claim-ok: our own ordering of the problems by class, each stated above

Problem~\ref{prob:open:zl} belongs to this class, because its local rings are
the local rings of a polynomial algebra.
Biswas, Gurjar and Kolte settle the affine surface case whenever the smooth
locus has logarithmic Kodaira dimension at most one
\cite[Theorem~1.2, p.~270]{BiswasGurjarKolte2014}.
Problem~\ref{prob:open:jc} lives here too, and it stays open for every $n\ge2$
\cite[pp.~1--2]{BelovKanelKontsevich2007}.
Problem~\ref{prob:open:dixmier} is the Weyl algebra form of the same question,
and Belov-Kanel and Kontsevich tie the two stable forms together
\cite[Theorem~1, p.~3]{BelovKanelKontsevich2007}.
Chapter~\ref{ch:polynomial-weyl} reads the Weyl algebra in
Section~\ref{sec:polynomial-weyl:weyl}.

The finite generation problems of Problems~\ref{prob:open:hdp}
and~\ref{prob:open:bass} also live in this class.
The Homogeneous Dependence Problem of Problem~\ref{prob:open:hdp} stays open
only in dimension $5$ \cite[p.~108]{Freudenburg2017}.
A homogeneous derivation of degree two with full rank is open in every
dimension \cite[Remark~3.51, p.~112]{Freudenburg2017}.
The triangularization problem of Problem~\ref{prob:open:bass} remains generally
open \cite[p.~288]{Freudenburg2017}.

The second class is the finite-dimensional algebra over a field.
Here the local property holds for some structures and fails for others.
Chapter~\ref{ch:local-characterizations} reads the local maps here in
Section~\ref{sec:local-characterizations:local}. % claim-ok: our own ordering of the problems by class, each stated above

Hadwin and Li settle the question of Larson for the $J$-subspace lattices: the
space $\operatorname{der}(\operatorname{alg}L)$ is reflexive
\cite[Theorem~2.15, p.~708]{HadwinLi2004}.
They also show that $\operatorname{der}(A,M)$ is algebraically $3$-reflexive for
every Banach $A$-bimodule $M$ \cite[Theorem~2.19, p.~710]{HadwinLi2004}.
The local property for automorphisms fails even in finite dimension, since
$T\mapsto AT^{t}A^{-1}$ on $M_n(\mathbb{C})$ is a surjective local automorphism
that is not an automorphism \cite[p.~713]{HadwinLi2004}.
So Problem~\ref{prob:open:larson} is answered for the lattices above, and the
general algebra of the question stays out of this class.

This class also holds a counterexample for the 2-local property: Zhang and Li
built a 2-local derivation on the upper triangular complex $2\times2$ matrices
that is not a derivation \cite[p.~15]{AyupovKudaybergenov2012}.
So the 2-local question of Problem~\ref{prob:open:twolocal} has no uniform answer
even among the small algebras of this class.

The third class is the $C^{*}$-algebra and its operator relatives.
Chapter~\ref{ch:local-characterizations} reads the 2-local notion in
Section~\ref{sec:local-characterizations:two-local}, and
Chapter~\ref{ch:banach-operator} reads the continuity question in
Section~\ref{sec:banach-operator:singer-wermer}. % claim-ok: our own ordering of the problems by class, each stated above

The 2-local problem of Problem~\ref{prob:open:twolocal} is settled here for the
largest algebras of the class.
Ayupov and Kudaybergenov prove that every 2-local derivation of $B(H)$ is a
derivation \cite[Theorem~2.3, p.~16]{AyupovKudaybergenov2012}.
Ayupov and Arzikulov extend this to every semi-finite von Neumann algebra
\cite[p.~9]{AyupovArzikulov2014}, and Ayupov and Kudaybergenov extend it to
every von Neumann algebra \cite[Theorem~2.1, p.~2]{AyupovKudaybergenov2015}.
Habibzadeh Fard and Sahleh add the compact operators
\cite[Corollary~3.4, p.~6]{HabibzadehFardSahleh2019}.
So the answer is yes for every algebra of the class that these papers reach.
The $\varphi$-derivation side of the class has one case that no one has settled
\cite[Remark~5.3, p.~720]{BresarVillena2002}.

\begin{problem}[{\cite[Remark~5.3, p.~720]{BresarVillena2002}}]\label{prob:class:phi-quasinilpotent}
Let $A$ be a unital Banach algebra, let $\varphi$ be an inner automorphism of
$A$, and let $\Delta$ be a $\varphi$-derivation of $A$.
If $a$ lies in the radical of $A$ and $\Delta^2(a)=0$, is $\Delta(a)$
quasinilpotent?
\end{problem}

Bre\v{s}ar and Villena prove Problem~\ref{prob:class:phi-quasinilpotent} for
derivations and report that they could not prove it for $\varphi$-derivations
\cite[Remark~5.3, p.~720]{BresarVillena2002}.
The continuity problem of Problem~\ref{prob:open:sw} is the other open question
of this class, and the authors call it one of the hardest problems of Banach
algebra theory \cite[p.~710]{BresarVillena2002}.

The fourth class is the Lie algebra, where the 2-local question has a negative
answer in general.
Chapter~\ref{ch:lie-algebras} reads the class in
Section~\ref{sec:lie-algebras:infinite-dimensional}. % claim-ok: our own ordering of the problems by class, each stated above

Ayupov and Yusupov answer Problem~\ref{prob:open:twolocal} with yes for the Witt
algebra \cite[Theorem~4.1, p.~4]{AyupovYusupov2020} and for the positive Witt
algebra \cite[Theorem~4.4, p.~6]{AyupovYusupov2020}.
They answer it with no for the thin Lie algebra, which admits a 2-local
derivation that is not a derivation \cite[Theorem~4.8, p.~8]{AyupovYusupov2020}.
So this class holds a counterexample of the book, and the problem for a
general infinite-dimensional Lie algebra stays open.

The fifth class is the prime ring, where the questions about the standard form
are answered for most characteristics.
Chapter~\ref{ch:lie-type} reads the Lie problem in
Section~\ref{sec:lie-type:martindale}, and Chapter~\ref{ch:higher-multi} reads
the symmetric bi-derivation in Section~\ref{sec:higher-multi:symmetric-maps}. % claim-ok: our own ordering of the problems by class, each stated above

Martindale proves the standard form for a primitive ring with a nontrivial
idempotent \cite[Theorem~2, p.~187]{Martindale1964}.
Bre\v{s}ar extends it to prime rings of characteristic not $2$ that do not
satisfy $S_4$ \cite[Theorem~5, p.~529]{Bresar1993LieTAMS}.
Lemma~1 of Bre\v{s}ar puts the rings that satisfy $S_4$ inside $M_2(F)$ over a
field $F$, up to the commutative case
\cite[Lemma~1, p.~530]{Bresar1993LieTAMS}.
So Problem~\ref{prob:open:herstein} stays open here for characteristic $2$ and
for a ring that satisfies $S_4$.
The symmetric bi-derivation of Problem~\ref{prob:open:vukman} stays open for a
power above $2$ \cite[pp.~188--189]{Vukman1990Two}, and the commuting map of
Problem~\ref{prob:open:commuting} stays open for a power above $1$
\cite[p.~334]{LeeTK1997Centralizing}.

The last class is the triangular algebra, and the problem here comes from the
$f$-derivation.
Chapter~\ref{ch:matrix-triangular} reads the class in
Section~\ref{sec:matrix-triangular:triangular}, and
Chapter~\ref{ch:generalized-twisted} reads the $f$-derivation side in
Section~\ref{sec:generalized-twisted:f-derivations}. % claim-ok: our own ordering of the problems by class, each stated above

Benkovi\v{c} proves the standard form for an $f$-derivation whose polynomial
has a nonzero sum of coefficients \cite[Theorem~1.1]{Benkovic2014}.
The remaining case is Problem~\ref{prob:open:fderiv}, and it stays open
\cite[Problem~1.2]{Benkovic2014}.

Across the six classes the same split appears. % claim-ok: our own summary of the six classes read above in this section
The polynomial class keeps the deepest problems open.
The operator classes settle the 2-local question for their large algebras, and
the small ones keep counterexamples: the thin Lie algebra and the upper
triangular complex $2\times2$ matrices.
Two problems have no class that settles them, and they are the reflexive space
of Larson and the noncommutative conjecture of Singer and Wermer.
They are the two that a later chapter is most likely to close first.
Each problem stated in the chapters appears here under the class where its
algebra lives.
The one problem that this section states itself sits with the
$C^{*}$-algebras.
So a change in one class can be read against the question family it belongs
to.
